\scanpage{180}

compacte --- le cas le plus simple de $X\simeq\mathbb P^1_{\mathbb C}\unsure{\simeq|\cong}S^2$, avec un diviseur lisse réduit à un point, donc $X^{*}\simeq\mathbb C$ (qui est bien un $K(\pi,1)$ avec $\pi=\{1\}$!), est le contre-exemple type: on ne reconstitue pas le type d'homotopie de $X=S^2$ \deleted{à partir de} \inserted{\unsure{comme}} étant \deleted{\unsure{celui}} \inserted{décrit} via le système ``\unsure{1-descriptif}'' de topos galoisien associé à

\diagram{\begin{tikzcd}[ampersand replacement=\&,column sep=large,row sep=small]
\& D_{\infty}^{*}\arrow[dl]\arrow[dr] \& \\
D_{\infty} \& \& \mathbb C\,\deleted{\unsure{\mathbb C^{*}}}=\mathbb P^1_{\mathbb C}\setminus\{\infty\}=X^{*}
\end{tikzcd}}

i.e. au diag.
\diagram{\begin{tikzcd}[ampersand replacement=\&,column sep=large,row sep=small]
\& \mathbb Z\arrow[dl]\arrow[dr] \& \\
\{\unsure{1}\} \& \& \{1\}
\end{tikzcd}}
des group\deleted{\unsure{es}}\inserted{\unsure{oïdes}}

\deleted{\unsure{qui se réduit par}} le topos global ``quasi-galoisien'' associé (faisceaux sur $S^2$ qui sont loc. constants donc constants sur $S^2\setminus\{\infty\}=\mathbb C$) s'identifie à $\widehat C$, où $C$ est la catégorie
\diagram{\begin{tikzcd}[ampersand replacement=\&]\bullet\arrow[r]\&\bullet\end{tikzcd}}
qui a un \unsure{objt} final, donc la cohomologie de $\widehat C$ est nulle en dim $>0$, alors que celle de $S^2$ ne l'est pas. Cela soulève la question générale de comprendre le type d'homotopie d'un topos obtenu par ``recollement'' d'un ouvert à son fermé complémentaire\ldots

\margin{179}
